\section{Discussion}
\label{sec:sn:discussion}

For the design spaces studied here, combinatorial growth and costly design-point
evaluation together make exhaustive exploration impractical; when the evaluation
dominates, reducing the number of evaluator calls is the main available leverage. The
results above quantify how much that leverage is worth. Reducing the expensive evaluations from the size of the
space to a few dozen, without giving up the guarantee that nothing optimal was discarded,
turns a study that would require a simulator campaign into one that fits inside a design
meeting; at the per-design-point costs reported for industrial trace-driven
exploration\cite{saadatmand2025compdse} the projection is two to three orders of
magnitude of wall-clock time.

The reason this works is narrower than ``pruning works'' and should be stated precisely.
Safe pruning of partial designs is an established idea in embedded system
synthesis\cite{neubauer2018exact,lukasiewycz2008symbolic} and in multi-objective search
more generally\cite{mandow2010namoa,paretopruning2022}. Its usual precondition ---
objectives that accumulate monotonically with each decision --- fails for a central objective in
latency-sensitive distributed cyber-physical systems, because latency depends on tier
utilisation and utilisation is fixed jointly by several decisions. The contribution here
is a way to keep exactness without that precondition: relax each additive term separately
over the domains still reachable, and order the decisions so that the contention term
stops being relaxed as early as possible. Three classes of objective term can be distinguished.
Terms that accumulate with each decision, such as deployment cost or transferred volume,
admit the usual bounds. Terms that depend on contention, such as queueing delay or
utilisation-dependent energy, admit exact bounds once the decisions that determine the
contention are fixed, and the traversal order is chosen so that this happens early. Terms
without a provable relaxation --- a black-box simulator output, for instance --- cannot be
bounded; they can still be evaluated on the survivors, but they cannot drive pruning.
The mechanism is not specific to the template
studied here. Formally, it applies whenever admissible term-wise relaxations can be constructed and a
subset of the architectural decisions determines the load-dependent terms. Empirical
transfer to a structurally different architecture template has not been demonstrated
here, so we make no empirical generality claim beyond the evaluated template. Deployment problems of that shape are common: placement and scaling
across a computing continuum\cite{filippini2026replica,predictive2026}, partitioned
inference pipelines whose accuracy and robustness depend on where each stage
runs\cite{guo2023,guo2025}, and energy-driven co-design of cloud--edge
serving\cite{cloud2026coopt}. Demonstrating that transfer is the natural next step.

The exactness provided here has a precise scope. The guarantee is that the method returns the exact Pareto front of the
analytical model; it says nothing about the model itself. A design-time method can be
exact with respect to its own model and still inherit every error of that model. This is
why the deployment round matters more than the correlation coefficient it produces: with
$25\%$ absolute error, the model is not a substitute for a simulator or a measurement,
but it orders clearly separated architectures the same way the hardware does in roughly
nine cases out of ten, and in \ReliabilityFifty{} cases once the predicted advantage exceeds
$50\%$. That monotone relation between predicted separation and decision reliability is
what a screening stage requires: it tells the engineer which comparisons can be settled on
the model and which must be sent to a simulator or a test bench. The productive
arrangement is therefore two-level. Analytical bounds remove the bulk of the space
exactly and cheaply; the few dozen survivors are where a high-fidelity simulator or a
physical deployment should spend its time. Nothing in the method requires the analytical
model to be the final word, and an evaluator-centric methodology plugged in behind it
would be strictly better than either component alone: trace-driven methodologies such as
CompDSE\cite{saadatmand2025compdse} supply the high-fidelity evaluator and leave the search
open, design-time tools for the computing continuum such as
SPACE4AI-D\cite{sedghani2024space4aid} search heuristically, and the method here supplies
the missing piece, a search that spends the evaluator only where optimality is still
possible and certifies what it skipped. The compact symbolic baseline makes the same point
from the other side. When the entire model is available in closed form, a factorised
constraint encoding also recovers the exact front, without calling any evaluator; the
reason to prune rather than solve is that the model which matters most is usually the one
that is not available in closed form.

Three limits apply to the results reported here. First, the size of the gain is set by the cost of one evaluation, not by the size of the
space. With the closed-form evaluator, which costs microseconds, the measured speed-up is
$\AnalyticalSpeedup\times$; the projected gains against the discrete-event evaluator are two
to three orders of magnitude larger. As the evaluator gets cheaper the search overhead
becomes the dominant term and the advantage narrows towards the ratio of overheads, so
the method should be chosen on evaluator cost rather than on space size. Second, the
guarantee is only as good as the bounds: a bound that cannot be proven admissible must
not be used, because an inadmissible bound silently converts an exact method into an
approximate one with no signal in the output. Third, near saturation the design-time
prediction breaks down and so does the deployment, as six of our measured configurations
show; architectures in that regime should be treated as unresolved rather than ranked.
Fourth, the deployment measured latency only. Energy and network volume are computed by
the model, and cost is a normalised relative tariff; the front is exact for all four
objectives of the model, but only its latency ordering has been checked against hardware.
Extending the guarantee to objectives that resist closed-form bounds, for example through
surrogate bounds that remain provably admissible, is the most valuable open direction we
see. More broadly, the results change what an analytical model is for in architecture
design: it need not replace a high-fidelity evaluator to be useful, because it can
certify which regions of a design space do not deserve one.
